\section{Conclusion}

This work extends hierarchical learning-guided exact Clifford+$T$ synthesis from full-register projective targets to explicit clean- and borrowed-ancilla contracts. The key semantic object is the isometry $U(C)J$ obtained by restricting the physical circuit to logical and borrowed inputs with every clean workspace qubit initialized in $\lvert0\rangle$. Correctness requires the requested logical transformation, identity on borrowed workspace, and exact restoration of clean workspace. Projective and exact-phase modes distinguish complete circuits from phase-sensitive modular subroutines.

The implementation preserves the original separation of responsibilities. A persistent DAG, forward--inverse Clifford tableau, ordered Pauli rotations, strengthened full-unitary canonicalizer, and resource-Pareto archive define the exact search. Outer linear SARSA allocates a frontier record; a role-aware disjoint LinUCB policy allocates a small batch of pending native continuations. Neither policy controls legality, equivalence, pruning validity, ancilla restoration, or certification.

The complete old and new test suites pass. Warm-started ancilla fine-tuning increases matched training time by approximately $8.2\%$. Three held-out clean-ancilla contracts are synthesized and certified. A 47-gate exact QFT-3 witness using one clean ancilla is independently verified with negligible numerical error and leakage, while a short unrestricted probe does not rediscover it. The resulting claim is therefore precise: the framework is functionally ancilla-aware and certifies nontrivial workspace-assisted Clifford+$T$ circuits, but long unrestricted ancilla synthesis remains a search challenge.

The next deterministic improvement is an exact contract-relative isometry canonicalizer, followed by larger compute--phase--uncompute benchmark families and a scalable independent verifier. These extensions can proceed without abandoning the mathematically transparent linear hierarchy established here.
